(a) Taking closures in $X$ carries a chain of irreducible closed subsets of $Y$ to a chain of the same length: if $Z$ is closed in $Y$, then $Y\cap\overline Z=Z$. Hence $\dim Y\le\dim X$.

(b) Given a chain $Z_0\subsetneq\cdots\subsetneq Z_n$ in $X$, choose $U_i$ meeting $Z_0$. Each $U_i\cap Z_j$ is irreducible and dense in $Z_j$ by \exref{I.1.6}, so the intersections are distinct and form a chain of length $n$. Thus $\dim X\le\sup_i\dim U_i$; the reverse inequality follows from (a).

(c) Take $X=\{0,1\}$ with open sets $\varnothing,\{1\},X$. Then $U=\{1\}$ is dense and open, but $\dim U=0<1=\dim X$.

(d) If $Y\ne X$, appending $X$ to a chain of length $\dim Y$ in $Y$ gives $\dim X\ge\dim Y+1$, a contradiction.

(e) On $\mathbb N$, take the closed sets to be $\varnothing$, $\mathbb N$, and the finite initial segments $\{1,\ldots,n\}$. Every descending chain stabilizes, but the initial segments are irreducible and give chains of arbitrary finite length.
